% !TeX root = ../main.tex
\section{The model and the algorithm}
\label{sec:methodology}

This section presents the model and the two samplers. Section~\ref{subsec:model} states the data model and the priors. Section~\ref{subsec:likelihood} introduces the working likelihood and its normal--exponential mixture representation, and Section~\ref{subsec:posterior} gives the posterior and relates it to quantile trend filtering. Section~\ref{subsec:individual} lists the steps of the sampler for a single quantile level, and Section~\ref{subsec:pmala} describes its trend update by P-MALA. Section~\ref{subsec:joint} extends the sampler to several quantile levels with ordered trends, and Section~\ref{subsec:joint_trend} gives its constrained trend and interweaving updates.

\subsection{Model and priors}
\label{subsec:model}

We observe $\by = (y_1, \ldots, y_n)^\top \in \R^n$ and model
\begin{equation}
    y_i = \theta_i + \varepsilon_i, \qquad i = 1, \ldots, n,
    \label{eq:datamodel}
\end{equation}
with independent errors, where for a level $\tau \in (0,1)$ the trend entry $\theta_i$ is the conditional $\tau$-quantile of $y_i$. Until Section~\ref{subsec:joint}, $\theta$, $z$, $\lambda$, and $\sigma$ refer to a single level $\tau$; Section~\ref{subsec:joint} attaches the level index.

Let $D = D^{(k+1)} \in \R^{m \times n}$ be the difference operator of order $k+1$, with $m = n - k - 1$ rows. Independent Laplace priors with common parameter $\lambda$ on the differences give
\begin{equation}
    p(\theta \mid \lambda) \propto \lambda^{m} \exp\!\left\{-\lambda \norm{D\theta}_1\right\}.
    \label{eq:prior_theta}
\end{equation}
The prior is flat on the null space of $D$, the polynomials of degree at most $k$, and $\lambda$, the penalty, controls the smoothness of the trend. The penalty and the likelihood scale $\sigma$ receive the conjugate priors
\begin{equation}
    \lambda \sim \mathrm{Gamma}(a_0, b_0), \qquad \sigma \sim \mathrm{IG}(a_\sigma, b_\sigma).
    \label{eq:hyperpriors}
\end{equation}
Following \citet{brantley2020baseline}, we use $k = 2$, so the prior favors piecewise quadratic trends.

\subsection{The likelihood and its mixture representation}
\label{subsec:likelihood}

The working likelihood is the asymmetric Laplace distribution (ALD) with scale $\sigma > 0$ and skewness $\tau$,
\begin{equation}
    f(u \mid \tau, \sigma) = \frac{\tau(1-\tau)}{\sigma}
    \exp\!\left\{-\frac{\rho_\tau(u)}{\sigma}\right\},
    \label{eq:ald}
\end{equation}
where $\rho_\tau(u) = u\{\tau - \mathbbm{1}(u < 0)\}$ is the check function \citep{koenker1978regression}. Its $\tau$-quantile is zero, and maximizing the likelihood in $\theta$ minimizes the check loss $\sum_i \rho_\tau(y_i-\theta_i)$. The check loss is not differentiable, so we use the normal--exponential representation of the ALD \citep{kozumi2011gibbs},
\begin{equation}
    \varepsilon_i = \psi_\tau z_i + \omega_\tau \sqrt{\sigma z_i}\, u_i, \qquad z_i \sim \mathrm{Exp}(\sigma), \qquad u_i \sim \mathcal{N}(0,1),
    \label{eq:ald_mixture}
\end{equation}
where $\mathrm{Exp}(\sigma)$ has mean $\sigma$, $z_i$ and $u_i$ are independent, $\psi_\tau = (1-2\tau)/\{\tau(1-\tau)\}$ and $\omega_\tau^2 = 2/\{\tau(1-\tau)\}$. Equivalently, $y_i \mid \theta_i, z_i, \sigma \sim \mathcal{N}(\theta_i + \psi_\tau z_i,\ \omega_\tau^2 \sigma z_i)$ with $z_i \mid \sigma \sim \mathrm{Exp}(\sigma)$. Integrating out $z_i$ returns \eqref{eq:ald}, so the representation leaves the model unchanged; the latent variables $z = (z_1,\ldots,z_n)^\top$ are sampled with the other parameters. Given $z$ the likelihood is Gaussian in $\theta$, with weights and pseudo-responses
\begin{equation}
    w_i = \frac{1}{\omega_\tau^2\, \sigma z_i}, \qquad r_i = y_i - \psi_\tau z_i, \qquad i = 1,\ldots,n.
    \label{eq:weights}
\end{equation}

\subsection{The posterior and quantile trend filtering}
\label{subsec:posterior}

The joint posterior of $(\theta, z, \lambda, \sigma)$ is
\begin{equation}
\begin{aligned}
    p(\theta, z, \lambda, \sigma \mid \by) \propto{} & \prod_{i=1}^n (\sigma z_i)^{-1/2} \exp\!\left\{-\frac{(y_i-\theta_i-\psi_\tau z_i)^2}{2\omega_\tau^2 \sigma z_i}\right\} \sigma^{-1} e^{-z_i/\sigma} \\
    & \times \lambda^{a_0 + m - 1} \exp\!\left\{-\lambda\left(b_0 + \norm{D\theta}_1\right)\right\} \sigma^{-a_\sigma - 1} e^{-b_\sigma/\sigma},
\end{aligned}
    \label{eq:posterior}
\end{equation}
and its marginal in $(\theta, \lambda, \sigma)$ is the posterior under the ALD likelihood.

With $z$ integrated out, the conditional posterior of $\theta$ at fixed $(\lambda, \sigma)$ is proportional to
\[
    \exp\Big\{-\sigma^{-1}\sum_{i=1}^n \rho_\tau(y_i-\theta_i) - \lambda\norm{D\theta}_1\Big\},
\]
with mode
\begin{equation}
    \hat{\theta}(\lambda, \sigma) \in \arg\min_{\theta \in \R^n} \left\{ \sum_{i=1}^n \rho_\tau(y_i-\theta_i) + \lambda\sigma \norm{D\theta}_1 \right\}.
    \label{eq:map_qtf}
\end{equation}
This is the quantile trend filtering problem of \citet{brantley2020baseline} with effective penalty $\lambda\sigma$. We sample $\lambda$ and $\sigma$ together with $\theta$ and report posterior means and pointwise intervals.

\subsection{The individual quantile sampler}
\label{subsec:individual}

We sample \eqref{eq:posterior} by a Gibbs sampler with four blocks, which we call Bayesian-PMALA. A sweep updates the blocks $\theta$, $z$, $\lambda$, and $\sigma$ in this order, each conditional on the most recent values of the others; the three conjugate conditionals are derived in Appendix~\ref{app:pmala:conditionals}. One sweep has five steps: the four block updates and an update of the P-MALA step size.
\begin{enumerate}
    \item \emph{Trend $\theta$.} Given $(z, \lambda, \sigma)$, the conditional density of $\theta$ is proportional to $\exp\{-U(\theta)\}$ with
    \begin{equation}
        U(\theta) = \frac{1}{2}\sum_{i=1}^n w_i (\theta_i - r_i)^2 + \lambda \norm{D\theta}_1 .
        \label{eq:Utheta}
    \end{equation}
    This conditional is not a standard distribution: $U$ is a weighted quadratic plus the $\ell_1$ penalty, convex but not differentiable wherever a difference $(D\theta)_j$ vanishes. We draw $\theta$ by P-MALA, a Metropolis--Hastings move built on the proximal map of $U$ (Section~\ref{subsec:pmala}).
    \item \emph{Latent variables $z$.} Given $(\theta, \sigma)$, the $z_i$ are independent with
    \begin{equation}
        z_i^{-1} \mid \theta, \sigma, \by \sim \mathrm{InvGauss}\!\left(\frac{1}{\tau(1-\tau)\,|y_i - \theta_i|},\ \frac{1}{2\tau(1-\tau)\,\sigma}\right),
        \label{eq:zupdate}
    \end{equation}
    where $\mathrm{InvGauss}(\mu, \nu)$ is the inverse Gaussian distribution with mean $\mu$ and shape $\nu$. The draw is exact, and it refreshes the weights $w_i$ and pseudo-responses $r_i$ that define $U$ in the next sweep.
    \item \emph{Penalty $\lambda$.} Given $\theta$,
    \begin{equation}
        \lambda \mid \theta, \by \sim \mathrm{Gamma}\!\left(a_0 + m,\ b_0 + \norm{D\theta}_1\right).
        \label{eq:lambdaupdate}
    \end{equation}
    Through this step the roughness $\norm{D\theta}_1$ of the current trend sets the amount of smoothing. Because posterior draws of the trend satisfy $\norm{D\theta}_1 \approx m/\lambda$, this update moves $\lambda$ slowly. It is therefore followed by an interweaving step \citep{yu2011center}, which rescales the component of $\theta$ outside the null space of $D$ together with $\lambda$ (Appendix~\ref{app:pmala:asis}).
    \item \emph{Likelihood scale $\sigma$.} Given $(\theta, z)$,
    \begin{equation}
        \sigma \mid \theta, z, \by \sim \mathrm{IG}\!\left(a_\sigma + \tfrac{3n}{2},\ b_\sigma + \sum_{i=1}^n \left\{ z_i + \frac{(y_i - \theta_i - \psi_\tau z_i)^2}{2\omega_\tau^2 z_i}\right\}\right).
        \label{eq:sigmaupdate}
    \end{equation}
    The scale enters the weights and the effective penalty $\lambda\sigma$ of \eqref{eq:map_qtf}, so it is drawn with the other parameters rather than fixed.
    \item \emph{Step size $\delta$.} The P-MALA step size is updated by a Robbins--Monro recursion toward acceptance rate $0.5$ (Appendix~\ref{app:pmala:adapt}).
\end{enumerate}
One run is stated in Algorithm~\ref{alg:pmala}. The point estimate is the posterior mean of $\theta$, and the $2.5\%$ and $97.5\%$ quantiles of the retained draws give pointwise $95\%$ intervals. Initialization, warm-up, and numerical safeguards are given in Appendices~\ref{app:pmala:init}--\ref{app:pmala:safeguards}.

\begin{algorithm}[h!]
\caption{Bayesian-PMALA for a single quantile level $\tau$}
\label{alg:pmala}
\begin{algorithmic}[1]
\Require Data $\by$, level $\tau$, order $k$, hyperparameters $a_0, b_0, a_\sigma, b_\sigma$, metric scale $\kappa$, sweeps $T_{\mathrm{gibbs}}$, burn-in $T_{\mathrm{burn}}$.
\State Initialize $\theta^{(0)} = \by$ and $(z^{(0)}, \lambda^{(0)}, \sigma^{(0)}, \delta^{(0)})$ \Comment{Appendix~\ref{app:pmala:init}}
\For{$t = 1, \ldots, T_{\mathrm{gibbs}}$}
    \State Form $w_i, r_i$ from \eqref{eq:weights} and $M = W + \kappa\lambda^2 D^\top D/2$; set $\gamma = \delta/2$
    \State Compute $T(\theta)$; draw $\theta^\star = T(\theta) + \sqrt{\delta}\,L^{-\top}\xi$; compute $T(\theta^\star)$ \Comment{Appendix~\ref{app:pmala:admm}}
    \State Accept $\theta^\star$ with probability $\min(1,\alpha)$, $\alpha$ from \eqref{eq:pmala_accept}
    \State Draw $z_i$ from \eqref{eq:zupdate}, $i = 1, \ldots, n$
    \State Draw $\lambda$ from \eqref{eq:lambdaupdate}; redraw $(\lambda, \theta)$ by the interweaving step
    \State Draw $\sigma$ from \eqref{eq:sigmaupdate}
    \State Update $\log\delta$ by a Robbins--Monro step toward acceptance $0.5$
\EndFor
\State \Return draws $\{\theta^{(t)}, \lambda^{(t)}, \sigma^{(t)}\}_{t > T_{\mathrm{burn}}}$
\end{algorithmic}
\end{algorithm}

\subsection{Sampling the trend by P-MALA}
\label{subsec:pmala}

The trend step requires a move that leaves the conditional of $\theta$ in \eqref{eq:Utheta} invariant. P-MALA \citep{pereyra2016proximal} replaces $U$ by its Moreau--Yosida envelope, draws a Langevin proposal from the envelope's gradient, and accepts or rejects it against the exact density. We apply it to the whole conditional \eqref{eq:Utheta}, which contains both the likelihood and the penalty.

The proposal is built from the proximal map of $U$. For a symmetric positive definite $M$ and $\gamma > 0$, the proximal map of $U$ in the $M$-norm and the envelope are
\begin{equation}
    T(x) = \arg\min_{\theta \in \R^n} \left\{ U(\theta) + \frac{1}{2\gamma}\norm{\theta - x}_M^2 \right\}, \qquad \norm{v}_M^2 = v^\top M v ,
    \label{eq:prox}
\end{equation}
\begin{equation}
    U^\gamma(x) = U\{T(x)\} + \frac{1}{2\gamma}\norm{T(x) - x}_M^2 .
    \label{eq:envelope}
\end{equation}
Because $U$ is convex, $T$ is unique and $U^\gamma$ is continuously differentiable with
\begin{equation}
    \nabla U^\gamma(x) = \frac{1}{\gamma}\, M\{x - T(x)\},
    \label{eq:envelope_grad}
\end{equation}
so its gradient costs one proximal map.

One Euler step of size $\delta$ of the Langevin diffusion for $\exp(-U^\gamma)$ preconditioned by $M^{-1}$ gives $\theta^\star = \theta - \frac{\delta}{2\gamma}\{\theta - T(\theta)\} + \sqrt{\delta}\, L^{-\top}\xi$, with $\xi \sim \mathcal{N}(\mathbf{0}, \mathbf{I}_n)$ and $M = L L^\top$. With $\gamma = \delta/2$ \citep{pereyra2016proximal} the proposal is
\begin{equation}
    \theta^\star \sim q(\cdot \mid \theta) = \mathcal{N}\!\left(T(\theta), \delta M^{-1}\right).
    \label{eq:pmala_proposal}
\end{equation}
The proposal is accepted with probability $\min(1,\alpha)$, where
\begin{equation}
    \log \alpha = U(\theta) - U(\theta^\star) + \frac{1}{2\delta}\left\{ \norm{\theta^\star - T(\theta)}_M^2 - \norm{\theta - T(\theta^\star)}_M^2 \right\}.
    \label{eq:pmala_accept}
\end{equation}
The move leaves the conditional \eqref{eq:Utheta} invariant: the acceptance ratio \eqref{eq:pmala_accept} uses $U$ itself and the Gaussian proposal density \eqref{eq:pmala_proposal}, and $T$ enters both only as a deterministic function of its input. The accuracy of the proximal solve therefore affects only the efficiency of the move.

The metric $M$ is chosen to match the curvature of $U$. We use
\begin{equation}
    M = W + \kappa\, \frac{\lambda^2}{2}\, D^\top D, \qquad W = \mathrm{diag}(w_1, \ldots, w_n), \qquad \kappa = 3,
    \label{eq:metric}
\end{equation}
which combines the precision of the augmented likelihood with a difference precision on the scale of the prior. The metric depends on $(z, \lambda, \sigma)$ but not on $\theta$, so the move is a valid Metropolis--Hastings kernel.

Problem \eqref{eq:prox} is a weighted trend filtering problem. We solve it by ADMM with a fixed number of iterations and a cold start, which makes $T$ a deterministic function of its input (Appendix~\ref{app:pmala:admm}). Every linear system is banded with half-bandwidth $k+1$, and one factorization serves the forward and reverse maps; the work per sweep is therefore linear in $n$. The step size $\delta$ is adapted toward acceptance $0.5$ (Appendix~\ref{app:pmala:adapt}).

\subsection{The joint quantile sampler}
\label{subsec:joint}

Quantile curves fitted one level at a time can cross. The joint sampler fits several levels with ordered trends. It is the sampler of Section~\ref{subsec:individual} applied level by level, except that each trend is drawn from its conditional restricted to lie between its neighbors, so its proximal map becomes a box-constrained problem. We first define the joint target and the bounds, then list the steps that differ from the individual sampler; Section~\ref{subsec:joint_trend} gives the constrained trend and interweaving updates.

Let $0 < \tau_1 < \cdots < \tau_J < 1$, and write $\theta_j$, $z_j$, $\lambda_j$, and $\sigma_j$ for the quantities of Sections~\ref{subsec:model}--\ref{subsec:individual} at level $\tau_j$. The levels are coupled only through the ordering
\begin{equation}
    \mathcal C = \left\{(\theta_1,\ldots,\theta_J):\
    \theta_{i1} \leq \cdots \leq \theta_{iJ},\quad i=1,\ldots,n\right\}.
    \label{eq:joint_order}
\end{equation}
With $\pi_j$ the augmented posterior \eqref{eq:posterior} at $\tau=\tau_j$, the joint working posterior is
\begin{equation}
    \pi_{\mathcal C}(\Theta,Z,\blambda,\bm\sigma\mid\by)
    \propto \mathbbm{1}_{\mathcal C}(\Theta)
    \prod_{j=1}^J \pi_j(\theta_j,z_j,\lambda_j,\sigma_j\mid\by),
    \label{eq:joint_target}
\end{equation}
where $\Theta=(\theta_1,\ldots,\theta_J)$ and $Z=(z_1,\ldots,z_J)$. Because the indicator depends only on $\Theta$, the conditionals of $z_j$, $\lambda_j$, and $\sigma_j$ are \eqref{eq:zupdate}--\eqref{eq:sigmaupdate} at $\tau = \tau_j$.

A sweep visits the levels in increasing order. At sweep $t$ the bounds of level $j$ are
\begin{equation}
\begin{aligned}
    \ell_j &= \begin{cases}\theta_{j-1}^{(t)}, & j>1,\\ -\infty\mathbf 1, & j=1,\end{cases}
    & u_j &= \begin{cases}\theta_{j+1}^{(t-1)}, & j<J,\\ +\infty\mathbf 1, & j=J,\end{cases} \\
    B_j &= \{\vartheta\in\R^n:\ \ell_j\leq\vartheta\leq u_j\},
\end{aligned}
    \label{eq:joint_box}
\end{equation}
held fixed during the trend and interweaving updates of level $j$, so the ordering holds after every accepted move. Given the other levels,
\begin{equation}
    p(\theta_j \mid \text{rest}) \propto \exp\{-U_j(\theta_j)\}\,\mathbbm{1}_{B_j}(\theta_j), \qquad
    U_j(\vartheta) = \frac12\sum_{i=1}^n w_{ij}(\vartheta_i-r_{ij})^2 + \lambda_j\norm{D\vartheta}_1,
    \label{eq:joint_conditional}
\end{equation}
the individual conditional \eqref{eq:Utheta} restricted to $B_j$.

Within a sweep, level $j$ is updated by the five steps of Section~\ref{subsec:individual}, with the following changes.
\begin{enumerate}
    \item \emph{Trend $\theta_j$.} The draw is the individual P-MALA move restricted to $B_j$: the proximal map is solved over the box, the metric $M_j$ and the step size $\delta_j$ belong to level $j$, and a proposal outside $B_j$ is rejected. This step and the interweaving step for $\lambda_j$ are the only ones that involve the other levels.
    \item \emph{Latent variables $z_j$.} Unchanged: \eqref{eq:zupdate} at $\tau=\tau_j$.
    \item \emph{Penalty $\lambda_j$.} The conjugate draw \eqref{eq:lambdaupdate} is unchanged. The interweaving step for $(\lambda_j,\theta_j)$ is restricted to the values for which $\theta_j$ stays in $B_j$.
    \item \emph{Likelihood scale $\sigma_j$.} Unchanged: \eqref{eq:sigmaupdate} at $\tau=\tau_j$.
    \item \emph{Step size $\delta_j$.} Each level has its own step size, updated toward acceptance $0.5$.
\end{enumerate}
During the first $H$ sweeps of the burn-in, the trend move is replaced by a constrained optimization step, and the interweaving step and the step-size update start afterwards (Appendix~\ref{app:pmala:init}). Algorithm~\ref{alg:joint_pmala} summarizes one run. Every retained draw satisfies \eqref{eq:joint_order}, so neither the draws nor the pointwise posterior means cross at the observed indices. For $J=1$ the joint sampler reduces to the individual sampler.

\begin{algorithm}[!htbp]
\caption{Joint Bayesian-PMALA for ordered quantile levels}
\label{alg:joint_pmala}
\begin{algorithmic}[1]
\Require Data $\by$, ordered levels $\tau_1,\ldots,\tau_J$, sweeps $T_{\mathrm{gibbs}}$, burn-in $T_{\mathrm{burn}}$, warm-up $H\leq T_{\mathrm{burn}}$.
\State Initialize ordered trends and each $(z_j,\lambda_j,\sigma_j,\delta_j)$ \Comment{Appendix~\ref{app:pmala:init}}
\For{$t=1,\ldots,T_{\mathrm{gibbs}}$}
    \For{$j=1,\ldots,J$}
        \State Fix $B_j$ from the latest neighbors \eqref{eq:joint_box}; form $w_{ij}, r_{ij}$, $M_j$; set $\gamma_j = \delta_j/2$
        \If{$t\leq H$}
            \State Replace $\theta_j$ by the constrained optimization warm-up step
        \Else
            \State Compute $T_j(\theta_j)$; draw $\theta_j^\star \sim \mathcal N(T_j(\theta_j), \delta_jM_j^{-1})$ \Comment{Appendix~\ref{app:pmala:boxsolver}}
            \State \textbf{if} $\theta_j^\star\notin B_j$ \textbf{then} reject; \textbf{else} compute $T_j(\theta_j^\star)$ and accept with probability $\min(1,\alpha_j)$, \eqref{eq:joint_accept}
        \EndIf
        \State Draw $z_j$ from \eqref{eq:zupdate} and $\lambda_j$ from \eqref{eq:lambdaupdate}
        \State \textbf{if} $t>H$ \textbf{then} redraw $(\lambda_j,\theta_j)$ by constrained interweaving in $B_j$
        \State Draw $\sigma_j$ from \eqref{eq:sigmaupdate}; update $\log\delta_j$ toward acceptance $0.5$ if $t > H$
    \EndFor
\EndFor
\State \Return draws $\{\theta_j^{(t)}, \lambda_j^{(t)}, \sigma_j^{(t)}\}_{t > T_{\mathrm{burn}},\ j = 1,\ldots,J}$
\end{algorithmic}
\end{algorithm}

\subsection{Constrained trend and interweaving updates}
\label{subsec:joint_trend}

Let $\iota_{B_j}$ be the indicator function of $B_j$ in the convex-analysis sense. The negative log-density $U_j + \iota_{B_j}$ is convex, so P-MALA applies as in Section~\ref{subsec:pmala}, with the metric $M_j$ of \eqref{eq:metric} at level $j$, a step size $\delta_j$ per level, and the box-constrained proximal map
\begin{equation}
    T_j(x) = \arg\min_{\vartheta\in B_j}
             \left\{U_j(\vartheta)+\frac{1}{2\gamma_j}\norm{\vartheta-x}_{M_j}^2\right\},
             \qquad \gamma_j=\delta_j/2 .
    \label{eq:joint_prox}
\end{equation}
The proposal is
\begin{equation}
    \theta_j^\star \sim q_j(\,\cdot\mid\theta_j) = \mathcal N\!\left(T_j(\theta_j),\ \delta_jM_j^{-1}\right).
    \label{eq:joint_proposal}
\end{equation}
A proposal outside $B_j$ has target density zero and is rejected. Otherwise it is accepted with probability $\min(1,\alpha_j)$,
\begin{equation}
    \log\alpha_j = U_j(\theta_j)-U_j(\theta_j^\star)
    +\frac{1}{2\delta_j}\left\{\norm{\theta_j^\star-T_j(\theta_j)}_{M_j}^2
            -\norm{\theta_j-T_j(\theta_j^\star)}_{M_j}^2\right\},
    \label{eq:joint_accept}
\end{equation}
and no truncation constant enters because the proposal itself is not truncated. When no bound binds, $T_j$ is the unconstrained map \eqref{eq:prox}, computed by the same ADMM; otherwise it is computed by a primal--dual interior-point method with banded Newton systems, so each map still costs $O(n)$ per iteration (Appendix~\ref{app:pmala:boxsolver}).

The interweaving step of the $\lambda_j$ block is applied within $B_j$. Write $p_j=QQ^\top\theta_j$ and $v_j=\theta_j-p_j$, with $Q$ an orthonormal basis of the null space of $D$. Holding $p_j$ and $\lambda_jD\theta_j$ fixed, a new value $\lambda_j'$ implies
\begin{equation}
    \theta_j(\lambda_j')=p_j+(\lambda_j/\lambda_j')\,v_j ,
    \label{eq:joint_asis_curve}
\end{equation}
and $\log\lambda_j'$ is drawn from the ancillary conditional restricted to the interval on which $\theta_j(\lambda_j')\in B_j$.
